\documentclass[fleqn,numbers=noenddot,headinclude,%twoside,%1headlines,%
				11pt,a4paper,footinclude,%
				cleardoublepage=empty,abstractoff %
                ]{scrartcl}

\usepackage[hyphens]{url}
\usepackage{float}

\def\tomas{Tom\'{a}\v{s} Peterka}
\def\tomasid{peterto8}
\def\franta{Franti\v{s}ek Hejl}
\def\frantaid{hejl}
\def\jan{J\'{a}n Tka\v{c}\'{i}k}
\def\janid{tkacijan}
\def\frantaid{hejl}

\def\laboratoryname{Ov\v{e}\v{r}en\'{i} metody CTEX}
\def\date{\today}


\input{preambl}

%\usepackage[subfigure]{tocloft}
%\renewcommand{\cftpartleader}{\cftdotfill{\cftdotsep}} % for parts
%\renewcommand{\cftchapleader}{\cftdotfill{\cftdotsep}} % for chapters
%\renewcommand{\cftsecleader}{\cftdotfill{\cftdotsep}} % for sections

\usepackage[sl,SF]{subfigure}
\usepackage{wrapfig}
\usepackage[hang,marginal]{footmisc}
\usepackage[justification=raggedright, margin=1cm,singlelinecheck=false,format=hang, labelfont={sl,small,rm},font={small,sf}, skip=10pt]{caption}
\renewcommand*\thesubfigure{\alph{subfigure}}

\DeclareMathOperator*{\argmax}{arg\,max}
\DeclareMathOperator*{\abs}{abs}

\makeatletter
\newcommand\wrapfill{\par
    \ifx\parshape\WF@fudgeparshape
    \nobreak
    \vskip-\baselineskip
    \vskip\c@WF@wrappedlines\baselineskip
    \allowbreak
    \WFclear
    \fi
}
\makeatother

\setlength{\intextsep}{6pt}

\usepackage{titlesec}
\titleformat{\section}{\large\bfseries}{\thesection}{1em}{}
\titleformat{\subsection}{\normalsize\bfseries}{\thesection}{1em}{}
\titleformat{\subsubsection}{\normalsize\bfseries}{\thesection}{1em}{}

\usepackage{amsmath}
\usepackage{graphicx}

\begin{document}

\makesavenoteenv[figure*]{figure}
\pagestyle{empty}
\onecolumn
%%Manuelle titelseite
\null  % Empty line
\nointerlineskip  % No skip for prev line
\vspace*{-1cm}
\begin{center}

\includegraphics[width=7cm,keepaspectratio=true]{./imgs/fit_logo.pdf}\\[0.3cm]
Czech Technical University\\
Faculty of Informatics

\vspace{1.2cm}

\Large MI-ROZ -- Rozpoznávání\\[0.5cm]
\laboratoryname \normalsize\\

\vspace{0.7cm}

\begin{tabular}{ccc}
\tomas & \franta & \jan \\
\tomasid @fit.cvut.cz & \frantaid @fit.cvut.cz & \janid @fit.cvut.cz \\
\end{tabular} \\[0.7cm]
\date
\end{center}
\vfill
\begin{abstract}
Práce popisuje implementaci algoritmu CTEX a jeho otestování na sérii zkušebních dat z ÚTIA (Ústav teorie informace a automatizace AV ČR).
\end{abstract}
\vfill
\tableofcontents

\clearpage

\pagestyle{scrheadings}
\setcounter{page}{1}

\section{Implementace CTex algoritmu}
\subsection{GB-FAB filter}
\label{sub:gb-fab}
Anisotropický filtr pro zvýraznění středně silných hran. Implementace podle zadání nebyla problém a nevznikly žádné odchylky. Jedině snad, že gradient je počítán 4 směrově místo 8. V článku není jasně definováno jak přesně jej počítat.
Tento filter používá upravenou hodnotu gradientu (pomocí mediánu z gradientu ve všech osách). Rozpuštění okolních pixelů do středového je řízena funkcí s průběhem na obrázku \ref{img:gbfab} a proto filter zvýrazňuje střední hrany, ponechává hrany ostré a zbytek celkem agresivně rozpouští do okolí.
\begin{figure}[htp]
\centering
\includegraphics[width=\textwidth]{plots/dfab_function.pdf}
\caption{Funkce rozpuštění gradientu do středu}
\label{img:gbfab}
\end{figure}
% subsection  (end)

\subsection{Výběr ideálního počtu shluků}
\label{sub:som}
Tato procedura sestává z více kroků. Nejdříve se pomocí SOM určí nejlepší počet centroidů a ty se poté předají do k-means algoritmu.

\subsubsection{Výběr dominantních barev} Slouží k iniciaci SOM mapy o velikosti 16 (čtverec $4\times4$), aby se nepoužívaly náhodné hodnoty. Prvník krokem je redukce barevného prostoru z (255, 255, 255) do (8, 8, 8) a zde sestrojíme 3D histogram. Z tohoto histogramu sebereme 16 barev s nejvyšší hodnotou pro iniciaci SOM mapy.
%subsubsection  (end)

\subsubsection{SOM analýza}
Po provedení a ukončení iterativního SOM algoritmu spočítáme váhy pro každý neuron v síti a spojujeme takové, které mají vzdálenost vah $<0.3$. Pokaždé když spojujeme 2 shluky, vybíráme jako reprezentanta ten s menší \textit{confidence} (vnitřní rozptyl částic). Výsledný počet neuronů/shluků je výstupem této části.
%subsubsection  (end)
% subsection  (end)

\subsection{Shlukování obrázků}
\label{sub:kmeans}
Shlukování obrázků je provedeno jak na RGB tak YIQ formátu. Počet shluků je sebrán z předchozího kroku \ref{sub:som}. Oba obrázky nejdříve projdou gb-fab filterem, poté jsou z nich extrahovány dominantní barvy a ty jsou, spolu s počtem shluků, použity jako iniciální nastavení k-means algoritmu. Výstupem je shluknutý obrázek jako například na obrázku \ref{img:clustered}, kde v levé části vidíme zdrojový a v pravé shluknutý.
\begin{figure}[htp]
\centering
\includegraphics[scale=1.00]{imgs/clustered.png}
\caption{Ukázka shluknutých obrázků}
\label{img:clustered}
\end{figure}
% subsection  (end)

\subsection{Multispace K-means}
\label{sub:multi-kmeans}
Zde je provedena druhá úroveň shlukování. Oba obrázky jsou spojeny, takže každý pixel je nyní reprezentován šesticí (R, G, B, Y, I, Q).
Centroidy jsou opět nastaveny podle dominantních barev (vypočítaných v minulém kroku) a také podle starého počtu shluků ze SOM analýzy.
% subsection  (end)

\subsection{Extrakce textur}
\label{sub:textury}
lorem ipsum
% subsection  (end)

\subsection{ASKM (Adaptive Spatial K-Means}
\label{sub:askm}
lorem ipsum
% subsection  (end)


\section{Vyhodnocení}
\label{sec:vyhodnoceni}
lorem ipsum
% section  (end)

\section{Závěr}
\label{sec:zaver}
lorem ipsum
% section  (end)
%\lstset{style=java}
%\begin{lstlisting}
%
%\end{lstlisting}

\end{document}
